\section{Conclusion}
\label{sec:conclusion}

This study evaluates a Tri-Domain ViT latent representation fusion framework (Face $\oplus$ Emotion $\oplus$ Age: 2,304 dimensions) for intersectional race and gender classification on the DemogPairs dataset. Experimental results demonstrate that tri-domain fusion achieves the best configuration across 3 of the 4 evaluated classifiers, encompassing SVM, LR, and GNB. Within the evaluated search space, the optimized SVM model with hyperparameter $C=10$ and a degree-two polynomial kernel yields the highest classification performance, namely an Accuracy of 93.70\% and an F1-Score of 93.69\%. Nevertheless, the benefit of multi-domain representation integration proves to be classifier-dependent, as RF achieves its peak performance on the dual-domain Emotion $\oplus$ Face scheme with an accuracy of 86.85\%. In addition, subgroup performance evaluation confirms that the F1-Score values across all six demographic subgroups are stably distributed within the range of 91.74\% to 96.14\%.

Despite these achievements, this study has several limitations that open avenues for future research. First, the demographic scope of this investigation is limited to three macro-racial groups (Asian, Black, and White) in the DemogPairs dataset, so the obtained empirical findings cannot be directly generalized to multiracial populations or other ethnic groups without additional comprehensive testing. Second, latent representation extraction is conducted separately (offline) utilizing frozen pre-trained ViT backbones without simultaneous end-to-end fine-tuning. Third, the experimental evaluation remains focused on controlled laboratory environment images, such that the model's response to extreme pose variations, occlusions, and unconstrained (in-the-wild) illumination requires further investigation. Therefore, future research directions are recommended to explore cross-modal adaptive attention fusion mechanisms (cross-attention transformer fusion) trained end-to-end to learn dynamic inter-domain interaction weights, as well as extending evaluation to larger, more heterogeneous datasets such as FairFace or UTKFace. In addition, the implementation of visual explainability methods such as Attention Rollout or Grad-CAM is crucial to transparently map the spatial contribution of each domain, accompanied by the enforcement of ethical artificial intelligence governance (Responsible AI) with human-in-the-loop oversight to prevent exploitation in public surveillance systems.
